% Options for packages loaded elsewhere
\PassOptionsToPackage{unicode}{hyperref}
\PassOptionsToPackage{hyphens}{url}
\PassOptionsToPackage{dvipsnames,svgnames,x11names}{xcolor}
\documentclass[
  10pt,
  fontset=fandol]{ctexart}
\usepackage{xcolor}
\usepackage[margin=16mm]{geometry}
\usepackage{amsmath,amssymb}
\setcounter{secnumdepth}{-\maxdimen} % remove section numbering
\usepackage{iftex}
\ifPDFTeX
  \usepackage[T1]{fontenc}
  \usepackage[utf8]{inputenc}
  \usepackage{textcomp} % provide euro and other symbols
\else % if luatex or xetex
  \usepackage{unicode-math} % this also loads fontspec
  \defaultfontfeatures{Scale=MatchLowercase}
  \defaultfontfeatures[\rmfamily]{Ligatures=TeX,Scale=1}
\fi
\usepackage{lmodern}
\ifPDFTeX\else
  % xetex/luatex font selection
\fi
% Use upquote if available, for straight quotes in verbatim environments
\IfFileExists{upquote.sty}{\usepackage{upquote}}{}
\IfFileExists{microtype.sty}{% use microtype if available
  \usepackage[]{microtype}
  \UseMicrotypeSet[protrusion]{basicmath} % disable protrusion for tt fonts
}{}
\makeatletter
\@ifundefined{KOMAClassName}{% if non-KOMA class
  \IfFileExists{parskip.sty}{%
    \usepackage{parskip}
  }{% else
    \setlength{\parindent}{0pt}
    \setlength{\parskip}{6pt plus 2pt minus 1pt}}
}{% if KOMA class
  \KOMAoptions{parskip=half}}
\makeatother
\setlength{\emergencystretch}{3em} % prevent overfull lines
\providecommand{\tightlist}{%
  \setlength{\itemsep}{0pt}\setlength{\parskip}{0pt}}
\usepackage{amsmath,amssymb,mathtools,needspace}
\xeCJKDeclareCharClass{CJK}{"2032,"0400->"04FF,"2160->"216B,"2460->"2473,"25A0,"25C7,"25CB,"25CF}
\IfFontExistsTF{Arial Unicode MS}{\xeCJKsetup{AutoFallBack=true}\setCJKfallbackfamilyfont{\CJKrmdefault}{Arial Unicode MS}}{}
\setcounter{secnumdepth}{0}
\setlength{\emergencystretch}{2em}
\allowdisplaybreaks
\usepackage{bookmark}
\IfFileExists{xurl.sty}{\usepackage{xurl}}{} % add URL line breaks if available
\urlstyle{same}
\hypersetup{
  pdftitle={误差的基本概念},
  colorlinks=true,
  linkcolor={Maroon},
  filecolor={Maroon},
  citecolor={Blue},
  urlcolor={Blue},
  pdfcreator={LaTeX via pandoc}}

\title{误差的基本概念}
\author{}
\date{}

\begin{document}
\maketitle

\subsubsection{1.3 误差的基本概念}\label{ux8befux5deeux7684ux57faux672cux6982ux5ff5}

\paragraph{1.3.1 误差与误差限}\label{ux8befux5deeux4e0eux8befux5deeux9650}

\textbf{定义1.1} 设 \(x\) 为准确值，\(x^*\) 为 \(x\) 的一个近似值，称 \(e^* = x^* - x\) 为近似值的绝

\begin{quote}
校注（agent 补充）：本页定义1.1末尾``绝''字与下一页``对误差''相接。
\end{quote}

\begin{quote}
校注（agent 补充，CH01-E001）：原书写``则 \(\tilde{I}_n\) 就是 \(E_0\) 的 \(n!\) 倍误差''，此处照录。由上式 \(E_n=(-1)^n n!E_0\) 可知，应理解为近似值 \(\tilde{I}_n\) 的误差绝对值为初值误差绝对值的 \(n!\) 倍；疑为原书表述不严谨，不是 OCR 错误。
\end{quote}

\href{../../source-images/pdf-018.jpeg}{查看原页}

对误差, 简称误差.

注意, 这样定义的误差 \(e^*\) 可正可负, 当绝对误差为正时近似值偏大, 叫做强近似值; 当绝对误差为负时近似值偏小, 叫做弱近似值.

通常, 我们不能算出准确值 \(x\), 也不能算出误差 \(e^*\) 的准确值, 只能根据测量工具或计算情况估计出误差的绝对值不超过某正数 \(\varepsilon^*\), 也就是误差绝对值的一个上界. \(\varepsilon^*\) 叫做近似值的误差限, 它总是正数. 例如, 用毫米刻度的米尺测量一长度 \(x\) (单位: mm, 下同), 读出和该长度接近的刻度 \(x^*\), \(x^*\) 是 \(x\) 的近似值, 它的误差限是 0.5, 于是 \(|x^* - x| \leqslant 0.5\); 如读出的长度为 765, 则有 \(|765 - x| \leqslant 0.5\). 从该不等式仍不知道准确的 \(x\) 是多少, 但知道 \(764.5 \leqslant x \leqslant 765.5\), 说明 \(x\) 在区间 \([764.5, 765.5]\) 上.

对于一般情形, \(|x^* - x| \leqslant \varepsilon^*\), 即 \(x^* - \varepsilon^* \leqslant x \leqslant x^* + \varepsilon^*\), 这个不等式有时也表示为

\[
x = x^* \pm \varepsilon^*.
\]

\paragraph{1.3.2 相对误差与相对误差限}\label{ux76f8ux5bf9ux8befux5deeux4e0eux76f8ux5bf9ux8befux5deeux9650}

误差限的大小还不能完全表示近似值的好坏. 例如, 有两个量 \(x = 10 \pm 1, y = 1000 \pm 5\), 则

\[
x^* = 10, \varepsilon_x^* = 1, y^* = 1000, \varepsilon_y^* = 5.
\]

虽然 \(\varepsilon_y^*\) 比 \(\varepsilon_x^*\) 大 4 倍, 但 \(\frac{\varepsilon_y^*}{y^*} = \frac{5}{1000} = 0.5\%\) 比 \(\frac{\varepsilon_x^*}{x^*} = \frac{1}{10} = 10\%\) 要小得多, 这说明 \(y^*\) 近似 \(y\) 的程度比 \(x^*\) 近似 \(x\) 的程度要好得多. 所以, 除考虑误差的大小外, 还应考虑准确值 \(x\) 本身的大小. 近似值的误差 \(e^*\) 与准确值 \(x\) 的比值

\[
\frac{e^*}{x} = \frac{x^* - x}{x}
\]

称为近似值 \(x^*\) 的相对误差, 记作 \(e_r^*\).

在实际计算中, 由于真值 \(x\) 总是不知道的, 通常取

\[
e_r^* = \frac{e^*}{x^*} = \frac{x^* - x}{x^*}
\]

作为 \(x^*\) 的相对误差, 条件是 \(e_r^* = \frac{e^*}{x^*}\) 较小, 此时

\[
\frac{e^*}{x} - \frac{e^*}{x^*} = \frac{e^* (x^* - x)}{x^* x} = \frac{(e^*)^2}{x^* (x^* - e^*)} = \frac{(e^* / x^*)^2}{1 - (e^* / x^*)}
\]

是 \(e_r^*\) 的二次方项级, 故可忽略不计.

相对误差也可正可负, 它的绝对值上界叫做相对误差限, 记作 \(\varepsilon_r^*, \varepsilon_r^* = \frac{\varepsilon^*}{|x^*|}\).

根据定义, \(\frac{\varepsilon_x^*}{|x^*|} = 10\%\) 与 \(\frac{\varepsilon_y^*}{|y^*|} = 0.5\%\) 分别为 \(x\) 与 \(y\) 的相对误差限, 可见 \(y^*\) 近似 \(y\) 的程度比 \(x^*\) 近似 \(x\) 的程度好.

\href{../../source-images/pdf-019.jpeg}{查看原页}

\paragraph{1.3.3 有效数字}\label{ux6709ux6548ux6570ux5b57}

当准确值 \(x\) 有多位数时，常常按四舍五入的原则得到 \(x\) 的前几位近似值 \(x^*\). 例如

\[
x = \pi = 3.14159265\cdots,
\]

取前三位，\(x_3^* = 3.14, \varepsilon_3^* \leqslant 0.002\)；取前五位，\(x_5^* = 3.1416, \varepsilon_5^* \leqslant 0.000008\)：它们的误差都不超过末位数字的半个单位，即

\[
|\pi - 3.14| \leqslant \frac{1}{2} \times 10^{-2}, \quad |\pi - 3.1416| \leqslant \frac{1}{2} \times 10^{-4}.
\]

若近似值 \(x^*\) 的误差限是某一位的半个单位，该位到 \(x^*\) 的第一位非零数字共有 \(n\) 位，就说 \(x^*\) 有 \(n\) 位有效数字. 如取 \(x^* = 3.14\) 作 \(\pi\) 的近似值，\(x^*\) 就有三位有效数字；取 \(x^* = 3.1416\) 作 \(\pi\) 的近似值，\(x^*\) 就有五位有效数字. \(x^*\) 有 \(n\) 位有效数字可写成标准形式

\[
x^* = \pm 10^m \times (a_1 + a_2 \times 10^{-1} + \cdots + a_n \times 10^{-(n-1)}), \tag{1.3.1}
\]

其中，\(a_1\) 是 1 到 9 中的一个数字；\(a_2, \cdots, a_n\) 是 0 到 9 中的一个数字；\(m\) 为整数，且

\[
|x - x^*| \leqslant \frac{1}{2} \times 10^{m-n+1}. \tag{1.3.2}
\]

\textbf{例1.2} 按四舍五入原则写出下列各数具有五位有效数字的近似数： 187.932 5，0.037 855 51，8.000 033，2.718 281 8.

\textbf{解} 按定义，上述各数具有五位有效数字的近似数分别是 187.93，0.037 856，8.000 0，2.718 3.

注意 \(x = 8.000033\) 的五位有效数字近似数是 8.0000 而不是 8，因为 8 只有一位有效数字.

\textbf{例1.3} 重力常数 \(g\)，如果以 \(\mathrm{m/s^2}\) 为单位，\(g \approx 9.80 \text{ m/s}^2\)；若以 \(\mathrm{km/s^2}\) 为单位，\(g \approx 0.00980 \text{ km/s}^2\)，它们都具有三位有效数字. 按第一种写法，有

\[
|g - 9.80| \leqslant \frac{1}{2} \times 10^{-2},
\]

据式(1.3.1)，这里 \(m=0, n=3\)；按第二种写法，有

\[
|g - 0.00980| \leqslant \frac{1}{2} \times 10^{-5},
\]

这里 \(m=-3, n=3\). 它们虽然写法不同，但都具有三位有效数字. 至于绝对误差限，由于单位不同，结果也不同，\(\varepsilon_1^* = \frac{1}{2} \times 10^{-2} \text{ m/s}^2,\allowbreak  \varepsilon_2^* = \frac{1}{2} \times 10^{-5} \text{ km/s}^2\)，而相对误差都是

\[
\varepsilon_r^* = 0.005/9.80 = 0.000005/0.00980.
\]

注意 相对误差与相对误差限是无量纲的，而绝对误差与误差限是有量纲的. 例 1.3 说明有效位数与小数点后有多少位数无关. 然而，从式(1.3.2)可得到具

\begin{quote}
校注（agent 补充）：本页末尾``具''字与下一页``有 \(n\) 位有效数字''相接。
\end{quote}

\href{../../source-images/pdf-020.jpeg}{查看原页}

有 \(n\) 位有效数字的近似数 \(x^*\), 其绝对误差限为

\[
\varepsilon^* = \frac{1}{2} \times 10^{m-n+1},
\]

在 \(m\) 相同的情况下, \(n\) 越大, \(10^{m-n+1}\) 越小, 故有效位数越多, 绝对误差限越小. 关于有效数字与相对误差限的关系, 有如下定理.

\textbf{定理1.1} 对于用式(1.3.1)表示的近似数 \(x^*\), 若 \(x^*\) 具有 \(n\) 位有效数字, 则其相对误差限为

\[
\varepsilon_r^* \leqslant \frac{1}{2a_1} \times 10^{-(n-1)};
\]

反之, 若 \(x^*\) 的相对误差限 \(\varepsilon_r^* \leqslant \frac{1}{2(a_1+1)} \times 10^{-(n-1)}\), 则 \(x^*\) 至少具有 \(n\) 位有效数字.

\textbf{证明} 由式(1.3.1)可得 \(a_1 \times 10^m \leqslant |x^*| \leqslant (a_1+1) \times 10^m\), 当 \(x^*\) 有 \(n\) 位有效数字时, 有

\[
\varepsilon_r^* = \frac{|x-x^*|}{|x^*|} \leqslant \frac{0.5 \times 10^{m-n+1}}{a_1 \times 10^m} = \frac{1}{2a_1} \times 10^{-n+1};
\]

反之, 有

\[
|x-x^*| = |x^*| \varepsilon_r^* \leqslant (a_1+1) \times 10^m \times \frac{1}{2(a_1+1)} \times 10^{-n+1} = 0.5 \times 10^{m-n+1},
\]

故 \(x^*\) 有 \(n\) 位有效数字. 证毕.

定理 1.1 说明, 有效位数越多, 相对误差限越小.

\textbf{例1.4} 要使 \(\sqrt{20}\) 的近似值的相对误差限小于 \(0.1\%\), 要取几位有效数字?

\textbf{解} 由定理 1.1 知

\[
\varepsilon_r^* \leqslant \frac{1}{2a_1} \times 10^{-n+1}.
\]

由 \(\sqrt{20} = 4.4 \dots\) 知, \(a_1 = 4\), 故只要取 \(n=4\), 就有

\[
\varepsilon_r^* \leqslant 0.125 \times 10^{-3} < 10^{-3} = 0.1\%,
\]

即只要对 \(\sqrt{20}\) 的近似值取四位有效数字, 其相对误差限就小于 \(0.1\%\), 此时

\[
\sqrt{20} \approx 4.472.
\]

\paragraph{1.3.4 数值运算的误差估计}\label{ux6570ux503cux8fd0ux7b97ux7684ux8befux5deeux4f30ux8ba1}

两个近似数 \(x_1^*\) 与 \(x_2^*\), 其误差限分别为 \(\varepsilon(x_1^*)\) 及 \(\varepsilon(x_2^*)\), 它们进行加、减、乘、除运算得到的误差限分别为

\[
\varepsilon(x_1^* \pm x_2^*) = \varepsilon(x_1^*) + \varepsilon(x_2^*),
\]

\[
\varepsilon(x_1^* x_2^*) \approx |x_1^*| \varepsilon(x_2^*) + |x_2^*| \varepsilon(x_1^*),
\]

\[
\varepsilon\left(\frac{x_1^*}{x_2^*}\right) \approx \frac{|x_1^*| \varepsilon(x_2^*) + |x_2^*| \varepsilon(x_1^*)}{|x_2^*|^2} \quad (x_2^* \neq 0).
\]

更一般的情况是, 当自变量有误差时, 计算函数值也会产生误差, 其误差限可利

\begin{quote}
校注（agent 补充）：本页末尾``可利''与下一页``用函数的 Taylor 展开式进行估计''相接。
\end{quote}

\href{../../source-images/pdf-021.jpeg}{查看原页}

用函数的 Taylor 展开式进行估计。设 \(f(x)\) 是一元函数，\(x\) 的近似值为 \(x^*\)，以 \(f(x^*)\) 近似 \(f(x)\)，其误差限记作 \(\varepsilon(f(x^*))\)，可用 Taylor 展开

\[
f(x)-f(x^*)=f'(x^*)(x-x^*)+\frac{f''(\xi)}{2}(x-x^*)^2,
\]

\(\xi\) 介于 \(x,x^*\) 之间。取绝对值得

\[
|f(x)-f(x^*)|\leqslant |f'(x^*)|\varepsilon(x^*)+\frac{|f''(\xi)|}{2}\varepsilon^2(x^*).
\]

假定 \(f'(x^*)\) 与 \(f''(x^*)\) 的比值不太大，可忽略 \(\varepsilon(x^*)\) 的高阶项，于是可得计算函数的误差限为

\[
\varepsilon(f(x^*))\approx |f'(x^*)|\varepsilon(x^*).
\]

当 \(f\) 为多元函数时计算 \(A=f(x_1,x_2,\cdots,x_n)\)，如果 \(x_1,x_2,\cdots,x_n\) 的近似值为 \(x_1^*,x_2^*,\cdots,x_n^*\)，则 \(A\) 的近似值为 \(A^*=f(x_1^*,x_2^*,\cdots,x_n^*)\)，于是函数值 \(A^*\) 的误差 \(e(A^*)\) 由 Taylor 展开，得

\[
\begin{aligned}
e(A^*)&=A^*-A=f(x_1^*,x_2^*,\cdots,x_n^*)-f(x_1,x_2,\cdots,x_n)\\
&\approx\sum_{k=1}^{n}\left(\frac{\partial f(x_1^*,x_2^*,\cdots,x_n^*)}{\partial x_k}\right)(x_k^*-x_k)
=\sum_{k=1}^{n}\left(\frac{\partial f}{\partial x_k}\right)^*e_k^*,
\end{aligned}
\]

于是误差限为

\[
\varepsilon(A^*)\approx\sum_{k=1}^{n}\left|\left(\frac{\partial f}{\partial x_k}\right)^*\right|\varepsilon(x_k^*);
\tag{1.3.3}
\]

而 \(A^*\) 的相对误差限为

\[
\varepsilon_r^*=\varepsilon_r(A^*)=\frac{\varepsilon(A^*)}{|A^*|}\approx\sum_{k=1}^{k}\left|\left(\frac{\partial f}{\partial x_k}\right)^*\right|\frac{\varepsilon(x_k^*)}{|A^*|}.
\tag{1.3.4}
\]

\textbf{例1.5} 已测得某场地长 \(l\) 的值为 \(l^*=110\ \mathrm{m}\)，宽 \(d\) 的值为 \(d^*=80\ \mathrm{m}\)，已知 \(|l-l^*|\leqslant 0.2\ \mathrm{m}\)，\(|d-d^*|\leqslant 0.1\ \mathrm{m}\)，试求面积 \(S=ld\) 的绝对误差限与相对误差限。

\textbf{解} 因 \(S=ld,\dfrac{\partial S}{\partial l}=d,\dfrac{\partial S}{\partial d}=l\)，由式(1.3.3)知

\[
\varepsilon(S^*)\approx\left|\left(\frac{\partial S}{\partial l}\right)^*\right|\varepsilon(l^*)+\left|\left(\frac{\partial S}{\partial d}\right)^*\right|\varepsilon(d^*),
\]

其中

\[
\left(\frac{\partial S}{\partial l}\right)^*=d^*=80\ \mathrm{m},\quad
\left(\frac{\partial S}{\partial d}\right)^*=l^*=110\ \mathrm{m},
\]

而

\[
\varepsilon(l^*)=0.2\ \mathrm{m},\quad\varepsilon(d^*)=0.1\ \mathrm{m},
\]

于是绝对误差限为

\[
\varepsilon(S^*)\approx(80\times0.2+110\times0.1)\ \mathrm{m}^2=27\ \mathrm{m}^2;
\]

相对误差限为

\[
\varepsilon_r(S^*)=\frac{\varepsilon(S^*)}{|S^*|}=\frac{\varepsilon(S^*)}{l^*d^*}\approx\frac{27}{8800}=0.31\%.
\]

\begin{quote}
校注（agent 补充，CH01-E002）：式(1.3.4)原图的求和上限清楚印为 \(k\)，本页保留 \(\sum_{k=1}^{k}\)。由函数有 \(n\) 个变量以及紧前式(1.3.3)的上限 \(n\) 判断，疑为原书将 \(n\) 误排为 \(k\)。OCR 候选写成 \(n\)，不能据此静默改原书。
\end{quote}

\begin{quote}
校注（agent 补充，CH01-E003）：原书关于可忽略高阶项的条件写为``\(f'(x^*)\) 与 \(f''(x^*)\) 的比值不太大''，此处照录。若意在使二阶项相对一阶项较小，应控制 \(|f''(\xi)|\varepsilon(x^*)/(2|f'(x^*)|)\)（在分母非零时）；原句疑有比值顺序或条件表述问题，不是 OCR 错误。
\end{quote}

\subsection{本节覆盖原页的校注记录}\label{ux672cux8282ux8986ux76d6ux539fux9875ux7684ux6821ux6ce8ux8bb0ux5f55}

下列记录按原页关联，含同页相邻小节的校注；计算时同时读取上文校注和完整原页。

\begin{itemize}
\tightlist
\item
  \texttt{CH01-E001}：原书 PDF 页 17
\item
  \texttt{CH01-E002}：原书 PDF 页 21
\item
  \texttt{CH01-E003}：原书 PDF 页 21
\end{itemize}

\end{document}
